% Generated draft: requires mathematical and source review.
\documentclass[11pt]{article}
\usepackage[margin=25mm]{geometry}
\usepackage{fontspec}
\setmainfont{Latin Modern Roman}
\usepackage{amsmath,amssymb,mathtools,unicode-math}
\setmathfont{Latin Modern Math}
\usepackage{enumitem,needspace,xurl,hyperref}
\setlength{\parindent}{0pt}
\setlength{\parskip}{5pt}
\setlength{\emergencystretch}{3em}
\newcommand{\N}{\mathbb{N}}
\newcommand{\Z}{\mathbb{Z}}
\newcommand{\Q}{\mathbb{Q}}
\newcommand{\R}{\mathbb{R}}
\newcommand{\C}{\mathbb{C}}
\newcommand{\F}{\mathbb{F}}
\newcommand{\A}{\mathbb{A}}
\newcommand{\PP}{\mathbb{P}}
\newcommand{\E}{\mathbb{E}}
\newcommand{\eps}{\varepsilon}
\newcommand{\dd}{\,\mathrm{d}}
\begin{document}
\section*{UnsolvedMath ID: 30004705}
\textbf{Source problem code:} catalogue entry
\subsection{Definitions and mathematical statement}
For a positive integer \(n\), a flat \(n\)-torus is a Riemannian manifold \(T_{\Lambda}=\R^n/\Lambda\), where \(\Lambda\subset\R^n\) is a full-rank lattice and the metric is induced by the Euclidean metric. Two such tori are isometric if there is a Riemannian isometry between them. Under the conventional interpretation used here, they are isospectral if their nonnegative Laplace--Beltrami operators on scalar functions have identical spectra, including multiplicities. More explicitly, writing
\[
\Lambda^*=\{\xi\in\R^n:\langle\xi,\lambda\rangle\in\Z\text{ for every }\lambda\in\Lambda\},
\]
the spectrum of \(T_{\Lambda}\) is the multiset \(\{4\pi^2\|\xi\|^2:\xi\in\Lambda^*\}\).

Using this interpretation, define the flat choir number by
\[
b_n=\sup\Bigl\{m\in\Z_{\geq 1}:\text{there exist flat }n\text{-tori }T_1,\ldots,T_m
\text{ such that, for all }1\leq i<j\leq m,
\ \operatorname{Spec}(T_i)=\operatorname{Spec}(T_j)
\text{ and }T_i\not\cong T_j\Bigr\},
\]
where \(\cong\) denotes Riemannian isometry and the supremum is taken in \(\Z_{\geq 1}\cup\{\infty\}\). The common spectrum is allowed to vary between collections.

\textbf{Problem.} Determine the exact value of \(b_n\) for every positive integer \(n\). In particular, determine \(b_{16}\), the largest possible cardinality of a collection of mutually isospectral, pairwise nonisometric flat \(16\)-tori.

\textbf{Unresolved source qualifications.} The supplied record does not explicitly define the spectral operator, the range of dimensions, or the term ``flat choir number.'' The definition above makes explicit the conventional scalar Laplace spectrum and the maximum-over-all-spectra interpretation suggested by the question. These conventions require source confirmation. The supplied dated assessment states that every \(b_n\) is finite and that new lower-bound constructions exist, while exact values in all dimensions, including dimension \(16\), remain undetermined. These are supplied curation claims, not independently verified mathematical or literature conclusions.
\subsection{Short English statement}
Determine, in each positive dimension, the maximum number of pairwise nonisometric flat tori sharing one scalar Laplace spectrum, with particular emphasis on dimension \(16\). The spectral and dimensional conventions require confirmation from the cited source; the supplied open status is a dated curation assessment.
\subsection{Sources}
\begin{itemize}
\item Local registry: \texttt{data/unsolvedmath/problems.json}, with \texttt{catalog/problems.json} as the supplementary registry; supplied permanent UnsolvedMath ID \(30004705\), source problem code \texttt{catalogue entry}, title ``Exact Flat Choir Numbers and Dimension-Sixteen Tori,'' in the Oberwolfach Reports Open Problems collection. The supplied record lists proposed year \(2021\) and does not identify an individual proposer. Registry files were not independently inspected for this draft.
\item Supplied report citation: \emph{Analysis, Geometry and Topology of Singular PDE} (2022), Oberwolfach Reports, pp.\ 1407--1461, DOI \texttt{10.4171/owr/2021/27}. \url{https://doi.org/10.4171/owr/2021/27}. The record describes this collection as extracted from workshop open-problem discussions. No problem-specific page or named speaker is supplied.
\item Supplied literature reference: Nilsson--Rowlett--Rydell, ``flat tori.'' \url{https://arxiv.org/abs/2110.09457}.
\item Supplied literature reference: M\"ardby--Rowlett--Rydell, ``six dimensions.'' \url{https://arxiv.org/abs/2412.16709}.
\item Literature assessment dated exactly \texttt{2026-08-22}: the supplied curation labels the problem open, reports finiteness of each flat choir number and new lower-bound constructions, and states that exact values in every dimension, including dimension \(16\), are not determined. This is supplied dated curation derived from the OWR record, not an independently verified assessment. The record itself requests independent verification of time-sensitive status and cited claims.
\item The supplied statement status is \texttt{corrected\_verified}, accompanied by a claim of manual semantic, typographic, and LaTeX review. This is source metadata and does not constitute independent mathematical review of this draft.
\end{itemize}
\end{document}
